\section{Introduction}

A vision-language model may correctly answer that the mug is left of the plate, that the spoon is closer than the fork, and that a ball moved behind a box. But if we interrogated the model the way a vision scientist interrogates an observer, would those answers add up to a single coherent percept of the world? This question matters because modern VLMs are increasingly asked to do more than recognize objects or describe images: they are evaluated on where objects are, how scenes are arranged, and what unfolds over time in videos. Yet success on question-answering benchmarks does not necessarily imply that a model has formed a stable spatial or temporal representation. A model may answer some questions correctly while relying on language priors, local appearance shortcuts, or mutually inconsistent heuristics. Recent work has made this tension increasingly visible: some benchmarks report progress on spatial reasoning, while others show that VLMs can ``see'' yet still fail to truly perceive, with responses heavily entangled with language priors \cite{chen2024spatialvlm,fu2024blink,chen2024mmstar,lee2025vlind}. Similar concerns have emerged in video understanding \cite{fu2024videomme,zhou2024mlvu,liu2024tempcompass,zhou2025vlm4d}. The problem, then, is not simply whether a model can answer spatial questions, but whether its answers reflect a coherent and visually grounded understanding of a world in space and time.

In biological vision research, progress rarely comes from a single aggregate score. It comes from controlled stimuli, dense behavioral probes, and carefully designed perturbations that reveal what an observer is sensitive to, what invariances it maintains, and what internal structure its behavior implies. Current VLM evaluation is often much coarser. Most benchmarks score isolated responses: they tell us whether a model selected the right option, but not whether its answers can be assembled into a single geometrically plausible scene, whether that scene is grounded in the image, or how that belief changes when visual evidence is removed, mismatched, or altered. Accuracy alone therefore cannot distinguish seeing from guessing. Two models may achieve similar question-level scores while exhibiting radically different behavior: one may express a coherent and visually grounded scene belief, while the other may produce a collection of locally plausible but globally incompatible judgments. What is missing is not just another benchmark, but an evaluation framework that treats a VLM as a black-box observer and recovers the structure implied by its behavior.

We address this problem with \textbf{Ordinary-Bench}, a controlled benchmark for probing VLM spatiotemporal understanding in rendered 3D environments. Inspired by the diagnostic role of synthetic benchmarks such as CLEVR and CLEVRER \cite{johnson2017clevr,yi2020clevrer}, and more broadly by the logic of psychophysics and systems neuroscience \cite{kriegeskorte2008rsa,yamins2016dicarlo}, Ordinary-Bench sacrifices ecological realism in exchange for exact geometric ground truth, controlled complexity, and repeatable perturbations. Our benchmark spans 700 procedurally generated scenes and more than 336{,}000 ordinal probes. Its static core consists of two task families: \emph{Quaternary Relative Relations} (QRR), which compare distances between object pairs, and \emph{Ternary Relative Relations} (TRR), which measure relative directional judgments. The framework also supports multi-view inputs and a temporal extension, but the central focus of this paper is a sharper question: when densely probed in a controlled world, what kind of scene belief does a VLM actually externalize?

To answer this question, Ordinary-Bench moves beyond question-level scoring and reconstructs the \emph{externalized scene belief} that best explains a model's responses. Model outputs are converted into spatial constraints, and the recovered scene is then evaluated for satisfiability, geometric fidelity, ambiguity, and distortion. This shift changes the unit of analysis. Rather than asking only how many local judgments are correct, we ask whether those judgments support a single coherent spatial explanation at all. Importantly, the goal is not to claim direct access to a model's internal neural state. Instead, we adopt a form of behavioral reverse engineering: infer the structure of a model's externalized scene belief from controlled probes, much as cognitive science uses structured tasks to study latent structure without equating behavior with mechanism. In this sense, Ordinary-Bench is both a benchmark and an evaluation methodology.

To separate genuine visual contribution from language priors, we complement this reconstruction framework with a three-condition protocol: correct image, mismatched image, and no image. Comparing behavior across these conditions allows us to estimate \emph{visual information gain} and to test whether apparent spatial competence is driven by visual grounding or by textual priors alone. We further incorporate selective answering and consistency analysis, following recent work showing that abstention, uncertainty, and response stability are essential for reliable multimodal evaluation \cite{whitehead2022reliablevqa,eisenschlos2024selectivevqa,khan2024consistency}. Together, these mechanisms yield a stricter notion of competence: a strong model should not only answer some spatial questions correctly, but should do so in a way that is reconstructable, internally consistent, and causally tied to the image.

Our contributions are three-fold:
\begin{itemize}
  \item We introduce \textbf{Ordinary-Bench}, a controlled benchmark for evaluating VLM spatial understanding in rendered 3D environments, with 700 scenes, over 336{,}000 ordinal probes, controlled complexity gradients, and support for multi-view and temporal extensions.
  \item We introduce \textbf{scene belief reconstruction} as a benchmark analysis tool: model responses are converted into spatial constraints and used to recover an externalized scene belief, enabling scene-level metrics such as constraint satisfaction, geometry preservation, and reconstruction error beyond question accuracy.
  \item We introduce a \textbf{visual contribution separation} protocol based on correct-image, mismatched-image, and no-image conditions, together with selective answering and consistency analyses, to quantify how much of a model's apparent spatial competence is genuinely grounded in vision.
\end{itemize}

Taken together, these components let us ask a sharper question than existing benchmarks typically permit: not merely whether VLMs can answer spatial questions, but what kind of world their behavior implies they actually see.
